% ===== PRÉAMBULE CANONIQUE — à coller VERBATIM en tête de CHAQUE .tex =====
\documentclass[11pt,a4paper]{article}
\usepackage[utf8]{inputenc}\usepackage[T1]{fontenc}\usepackage{lmodern}
\usepackage[french]{babel}
\usepackage{amsmath,amssymb,mathtools}\usepackage{icomma}
\usepackage[version=4]{mhchem}          % \ce{...} : équations & formules
\usepackage{siunitx}\sisetup{locale=FR} % \qty{}{}, \si{}, \ang{}
\usepackage{chemfig}                    % \chemfig{...} : structures développées
\usepackage{xcolor}\usepackage{colortbl}
\usepackage{tikz}\usepackage{pgfplots}\pgfplotsset{compat=1.18}
\usepackage[most]{tcolorbox}\usepackage{enumitem}\usepackage{array,booktabs}
\usepackage[a4paper,margin=2.2cm]{geometry}
\usetikzlibrary{arrows.meta,calc,angles,quotes,positioning,patterns,backgrounds,shapes.geometric}
\definecolor{primary}{RGB}{222,49,99}\definecolor{primarydark}{RGB}{150,22,55}
\definecolor{defcol}{RGB}{0,114,178}\definecolor{methcol}{RGB}{52,92,148}
\definecolor{excol}{RGB}{0,135,90}\definecolor{attncol}{RGB}{200,40,55}
\tcbset{boxbase/.style={enhanced,breakable,boxrule=0.9pt,arc=2.5pt,
  left=3mm,right=3mm,top=2.3mm,bottom=2.3mm,fonttitle=\bfseries,coltitle=white}}
% --- boîtes TUTORIEL (noms EXACTS, ne pas renommer) ---
\newtcolorbox{cle}[1][]{boxbase,colback=defcol!6,colframe=defcol,
  title={\textsc{À retenir}\ifx\relax#1\relax\relax\else\textmd{ : \itshape #1}\fi}}
\newtcolorbox{methode}[1][]{boxbase,colback=methcol!7,colframe=methcol,sharp corners,
  title={\textsc{Méthode}\ifx\relax#1\relax\relax\else\textmd{ --- \itshape #1}\fi}}
\newtcolorbox{exemplebox}[1][]{boxbase,colback=excol!5,colframe=excol,
  title={\textsc{Exemple}\ifx\relax#1\relax\relax\else\textmd{ : \itshape #1}\fi}}
\newtcolorbox{piege}[1][]{boxbase,colback=attncol!6,colframe=attncol,
  title={\textsc{Piège}\ifx\relax#1\relax\relax\else\textmd{ : \itshape #1}\fi}}
\newtcolorbox{remarque}[1][]{boxbase,colback=black!4,colframe=black!45,
  title={\textsc{Remarque}\ifx\relax#1\relax\relax\else\textmd{ : \itshape #1}\fi}}
% --- boîtes EXERCICE / CORRIGÉ (noms EXACTS, auto-comptées) ---
\newtcolorbox[auto counter]{exo}[1][]{boxbase,colback=primary!4,colframe=primary,
  title={\textsc{Exercice}~\thetcbcounter\ifx\relax#1\relax\relax\else\textmd{ --- \itshape #1}\fi}}
\newtcolorbox[auto counter]{corrige}[1][]{boxbase,colback=excol!4,colframe=excol,
  title={\textsc{Corrigé}~\thetcbcounter\ifx\relax#1\relax\relax\else\textmd{ --- \itshape #1}\fi}}
\newcommand{\R}{\mathbb{R}}\newcommand{\N}{\mathbb{N}}\newcommand{\Z}{\mathbb{Z}}\newcommand{\ds}{\displaystyle}
% (un \usepackage{fancyhdr}+en-tête est possible ; il est ignoré côté web)
% =========================================================================
\begin{document}

\begin{center}{\large\bfseries\color{primarydark} Thème 1 — Niveau Difficile}\end{center}
\emph{Données (\si{\gram\per\mole}) :} $\ce{H}=1$, $\ce{C}=12$, $\ce{N}=14$, $\ce{O}=16$,
$\ce{S}=32$, $\ce{K}=39$, $\ce{Cr}=52$, $\ce{Cu}=63{,}5$. \quad $N_{\mathrm{A}} = 6{,}02\times10^{23}\ \si{\per\mole}$.

\begin{exo}[retrouver une formule à partir des pourcentages massiques]
Un composé ne contient que du potassium, du chrome et de l'oxygène. Son analyse donne
$26{,}6\,\%$ de potassium et $35{,}4\,\%$ de chrome en masse (le reste étant de l'oxygène).
Détermine sa formule brute \ce{K_xCr_yO_z}.
\emph{Indication : raisonne sur $\SI{100}{\gram}$ de composé.}
\end{exo}

\begin{exo}[égaliser des nombres d'atomes]
Quelle masse de dioxyde de carbone \ce{CO2} contient \emph{autant} d'atomes d'oxygène que
$\SI{3.6}{\gram}$ d'eau \ce{H2O}~? ($M(\ce{H2O})=\SI{18}{\gram\per\mole}$, $M(\ce{CO2})=\SI{44}{\gram\per\mole}$.)
\end{exo}

\begin{exo}[identifier un métal par sa masse molaire]
On pèse un échantillon d'un métal pur : $\SI{0.15}{\mole}$ de ce métal a une masse de $\SI{9.525}{\gram}$.
\begin{enumerate}[label=\alph*)]
  \item Détermine sa masse molaire $M$.
  \item De quel métal peut-il s'agir (parmi ceux des données)~?
  \item Combien d'atomes contient l'échantillon~?
\end{enumerate}
\end{exo}

\begin{exo}[remonter à la formule d'un oxyde de soufre]
Un flacon contient $6{,}02\times10^{22}$ molécules d'un oxyde de soufre \ce{SO_x}, de masse totale
$\SI{8}{\gram}$.
\begin{enumerate}[label=\alph*)]
  \item Quelle est la quantité de matière $n$ présente~?
  \item En déduire la masse molaire $M$ du composé.
  \item Déterminer $x$ et nommer l'oxyde.
\end{enumerate}
\end{exo}

\begin{exo}[bilan d'un mélange]
Un mélange gazeux contient $\SI{0.2}{\mole}$ de diazote \ce{N2} et $\SI{0.3}{\mole}$ de
dioxygène \ce{O2}. ($M(\ce{N2})=\SI{28}{\gram\per\mole}$, $M(\ce{O2})=\SI{32}{\gram\per\mole}$.)
\begin{enumerate}[label=\alph*)]
  \item Quelle est la masse totale du mélange~?
  \item Quel est le nombre total de molécules~?
  \item Quel est le nombre total d'atomes~?
\end{enumerate}
\end{exo}

\end{document}
